\section{Limitations}
\label{sec:limitations}

\subsection{Evidence status}

This manuscript reports no trained research checkpoint and no benchmark result.
The implementation has been smoke-tested, but model capability, safety, and
efficiency claims remain untested. Readers should treat the experimental
sections as a protocol, not a finding.

\subsection{Model capacity}

A 0.8B-parameter model has limited parametric knowledge and limited multi-step
reasoning. It may memorize common facts while failing on novel combinations,
version-specific advice, or long contexts. Fine-tuning on cybersecurity data can
improve domain style and recall but cannot reliably create capabilities absent
from the base model. The project therefore includes retrieval and explicit
uncertainty rather than assuming that specialization solves knowledge limits.

\subsection{Architecture and tooling}

Qwen3.5's hybrid linear/full attention and multimodal checkpoint require a recent
Transformers implementation. Compatibility can change across library releases.
The CPU reference path is correct but slow; optional accelerated kernels are not
installed by default. The vision tower is unused, and the multimodal capability
of the checkpoint is outside scope.

\subsection{Data and licensing}

Public cybersecurity sources have different licenses and quality. Some are
copyleft or share-alike; some documentation is proprietary. Derived datasets and
models may inherit obligations. The project records licenses but cannot provide
legal clearance. NVD and ATT\&CK descriptions can be incomplete or outdated.
Teacher-generated data can contain errors and may be constrained by the teacher's
license. The committed smoke data is synthetic and unrepresentative.

\subsection{Contamination}

The 13-gram screen catches verbatim overlap only. Paraphrased benchmark leakage,
memorization, and semantic contamination may remain. Benchmark test items must be
excluded by construction, and the paper must report versions and subset choices.
No automated screen can prove absence of contamination.

\subsection{Evaluation validity}

Multiple-choice accuracy, token F1, and lexical grounding proxies are imperfect.
A model can score well by learning answer formats or by exploiting artifacts.
False-positive rates depend on the chosen positive label and class balance.
Syntax validity is not functional correctness, and the harness does not execute
generated code. Human review is required for safety and factual claims, and
inter-annotator agreement should be reported when human evaluation is used.

\subsection{Hardware and cost}

Memory estimates are approximations. Actual peak memory depends on kernels,
allocator behavior, sequence length, batch size, and concurrent workloads.
Consumer hardware varies widely. Cost is only meaningful with explicit rates and
a stated utilization assumption. Single-machine latency numbers are not portable.

\subsection{Deployment and safety}

The local API is a compatibility subset, not a production serving system. It has
no sandbox for generated commands because it does not execute them; integrators
who add execution own the sandbox and authorization controls. Prompt injection
through retrieved documents or code remains a risk. Small models can be
confidently wrong, and security decisions require human verification.

\subsection{Generalization}

The planned corpus emphasizes English-language, Linux/networking, and
web-security content. Results may not transfer to Windows-heavy enterprises,
industrial control systems, non-English environments, or proprietary tooling.
Larger models and additional data are needed to test generality.
